\documentclass[10pt]{article}

\usepackage[utf8]{inputenc}

\usepackage[T1]{fontenc}

\usepackage{amsmath}

\usepackage{amsfonts}

\usepackage{amssymb}

\usepackage[version=4]{mhchem}

\usepackage{stmaryrd}

\usepackage{bbold}

\usepackage{mathrsfs}



\begin{document}

сингулярностей, ведет переход к полярным координатам $\rho(x, \tau), \varphi(x, \tau)$ медленного поля



\[

\psi_{0}(x, \tau)=(\rho(x, \tau))^{1 / 2} \exp (i \varphi(x, \tau)) .

\]



Во многих случаях к цели ведет переход в $\Phi$. и. по медленным полям к переменным, соответствующим коллективным степеням свободы системы.



В теориях типа электродинамики, в к-рых массивные ферми- или бозе-частицы взаимодействуют с безмассовым векторным полем $A_{\mu}$, полезна формула



\[

G(x, y \mid A) \cong G(x, y \mid 0) \exp \left(i e \int_{x}^{y} A_{\mu} d x^{\mu}\right)

\]



для функции Грина оператора Дирака (или Д'Аламбера) в медленном поле $A_{\mu}$, где $\int_{x}^{y} \mathrm{~A}_{\mu} d x^{\mu}$ - интеграл по прямой, соединяющей точки $x$ и $y$ пространства-времени. Функциональное интегрирование выражений типа (5) по медленным полям дает инфракрасную асимптотику функций Грина в квантовой электродинамике, формулы дважды логарифмич. асимптотики, различные модификации эйкональных формул.



Jum.: [1] Wi en e r N., «J. Math. Phys.», 1923, v. 2, № 3, p. 131-74; [2] е го же, «Proc. London Math. Soc.», 1924, v. 22, № 6, p. 454-67; [3] Фей н ман Р., Хи бс А., Квантовая механика и интегралы по траекториям, пер. с англ., М., 1968; [4] Гельфанд И. М., Ягло м А. М., «Успехи математич. наук», 1956, т. 11, в. 1, с. 77-114; [5] Бо гол юбов Н. Н., Ш и рков Д. В., Введение в теорию квантованных полей, 4 изд., М., 1984; [6] Ковальч и к И. М., Ян ов ич Л.А., Обобщенный винеровский интеграл и некоторые его приложения, Минск, 1989; [7] Фадде е в Л.Д., «Теоретич. и математич. физика», 1969, т. 1, № 1, с. 3-17; [8] П о по в В. Н., Континуальные интегралы в квантовой теории поля и статистической физике, М., 1976; [9] Po pov V. N., Functional integrals and collective excitations, Camb., 1988; [10] По пов В. Н., Федотов С. А., «Журнал эксперим. и теоретич. физики», 1988, т. 94, в. 3, с. 183-94.



В. H. Попов.



ФУPЬE ГИПЕРФУНЦИЯ - модификация г гиерфункции, тесно связанная с аналитическими функционалами и обладающая преобразованием Фурье. Ф. г. были введены М. Сато (M. Sato) одновременно с гиперфункциями в конце 50-х гг. 20 в. Теория Ф. г. построена Т. Каваи [1]. Ф. г. широко применяются в теории линейных уравнений с частными производными. Имеются приложения Ф. г. в аксиоматич. квантовой теории поля.



Пусть $\mathbb{D}^{n}=\mathbb{R}^{n} \cup S^{n-1}$ - компактификация $\mathbb{R}^{n}$ с помощью сферы «бесконечного» радиуса. Для всякого открытого подмножества $U \subset \mathbb{D}^{n}+i \mathbb{R}^{n}$ через $x_{\varepsilon}(U)$ обозначено банахово пространство всех функций $\Phi(z)$, аналитических в $U \cap \mathbb{C}^{n}$ с кеонечной нормой



\[

|\Phi|=\sup \left\{e^{\varepsilon|z|}|\Phi(z)| ; z \in U \cap \mathbb{C}^{n}\right\}

\]



Пусть $\Omega$ - открытое подмножество из $\mathbb{D}^{n}+i \mathbb{R}^{n}$, и



\[

\begin{aligned}

& x_{-}(\Omega)=\underset{U \Subset \Omega}{\lim \operatorname{proj} \lim \operatorname{proj} x_{\varepsilon}(U),} \\

& x_{+}(\Omega)=\operatorname{limproj}_{U \Subset \Omega} \operatorname{limind}_{\varepsilon>0} x_{\varepsilon}(U),

\end{aligned}

\]



запись $U \Subset \Omega$ означает, что замыкание $\bar{U}-$ компакт в $\mathbb{D}^{n}+i \mathbb{R}^{n}$ и $\bar{U} \subset \Omega$. Для компакта $K \subset \mathbb{D}^{n}+i \mathbb{R}^{n}$ пусть



\[

x_{+}(K)=\operatorname{limind}_{U \supset K} \operatorname{limind}_{\varepsilon>0} x_{\varepsilon}(U) .

\]



Пространства $x_{-}(\Omega)$ вместе с естественными сужениями $x_{-}(\Omega) \rightarrow x_{-}(\omega)$ при $\Omega \subset \omega$ образуют пучок $x_{-}$ростков аналитич. функщий в $\mathbb{D}^{n}+i \mathbb{R}^{n}$, растущих в вешественных направлениях медленнее любой линейной экспоненты.



По аналогии с гиперфункциями пространство $\mathscr{R}(O) \Phi$. г. в области $O \subset \mathbb{D}^{n}$, по определению, есть $H_{O}\left(\Omega ; x_{-}\right)-n$-я группа когомологий множества $\Omega$ с коэффициентами в пучке $x_{-}$и с носителями в $O$, где $\Omega-$ произвольная окрестность $O$ в $\mathbb{D}^{n}+i \mathbb{R}^{n}$.



Ф. г. обладают всеми свойствами, присущими гиперфункциям. Каждая Ф. г. $g \in \Re(O)$ задается конечным набором $\left\{f_{j}\right\}$ функций $f_{j} \in x_{-}\left(D_{j}(O)\right)$, где $D_{j}(O)-$ тубоиды в $\mathbb{D}^{n}+i \mathbb{R}^{n}$, и обратно, всякий такой набор задает Ф. г. Нулевые Ф. г. представляются функциями



\[

f_{j}(z)=\sum f_{j k}(z), \quad z \in D_{j}(O),

\]



где функции $f_{j k}=-f_{k j} \in x_{-}\left(D_{j k}(O)\right)$, тубоиды $D_{j k}(O)=$ $=D_{k j}(O) \supset D_{j}(O) \cup D_{k}(O)$. Как и в случае гиперфункций, для подобласти $\omega$ области $O$ определяется сужение $\Re(O) \rightarrow$ $\rightarrow \mathfrak{A}(\omega)$, и Ф.г. образуют вялый пучок $\Re$ линейных пространств на $\mathbb{D}^{n}$.



Сужение пучка Ф. г. $\mathscr{R}$ с $\mathbb{D}^{n}$ на $\mathbb{R}^{n}$ совпадает с пучком гиперфункций $\mathscr{B}$.



Множество Ф. г. $g \in \Re\left(\mathbb{D}^{n}\right)$, носители к-рых лежат в компакте $K \subset \mathbb{D}^{n}$, отождествляется с сопряженным пространством $x_{+}^{\prime}(K)$. В отличие от $\mathbb{R}^{n}$, каждое замкнутое подмножество $\mathbb{D}^{n}$ есть компакт. В частности, $\mathcal{A}\left(\mathbb{D}^{n}\right)=x_{+}^{\prime}\left(\mathbb{D}^{n}\right)$, то есть каждая Ф. г. представляется аналитич. функционалом.



Основное пространство



\[

x_{+}\left(\mathbb{D}^{n}\right)=\operatorname{limind}_{\varepsilon>0} x_{\varepsilon}\left(\mathbb{R}^{n}+i B_{\varepsilon}\right),

\]



где шар $B_{\varepsilon}=\{|y|<\varepsilon\}$. На $x_{+}\left(\mathbb{D}^{n}\right)$ действует преобразование Фypьe $F$ :



\[

x_{+}\left(\mathbb{D}^{n}\right) \cong x_{+}\left(\mathbb{D}^{n}\right),

\]



так что определено преобразование Фурье Ф. г.:



\[

F: \Re\left(\mathbb{D}^{n}\right) \cong R\left(\mathbb{D}^{n}\right)

\]



(см. Фурье преобразование обобщенной функции).



На пространствах Ф. г. действуют дифференциальные операторы бесконечного порядка, такие же, как и на гиперфункциях.



Jum.: [1] Kaw a i T., «J. Fac. Sci. Univ. Tokyo, Sect. IA Math.», 1970, v. 17, p. 467-517. B. В. Жаринов. ФУPЬE ИНTEГPAЛ - формула для разложения непериодичеєкой функции на гармонические компоненты, частоты которых пробегают непрерывную совокупность значений. Если функция $f(x)$ абсолютно интегрируема на каждом конечном отрезке и если сходится



\[

\int_{-\infty}^{\infty}|f(x)| d x

\]



To



\[

f(x)=\frac{1}{\pi} \int_{0}^{\infty} d u \int_{-\infty}^{\infty} f(t) \cos u(x-t) d t .

\]



Эта формула впервые встречается при решении нек-рых задач теплопроводности у Ж. Фурье (J. Fourier, 1811), но еe доказательство было дано позже другими математиками. Формулу (1) можно представить также в виде:



\[

f(x)=\int_{0}^{\infty}[a(u) \cos u x+b(u) \sin u x] d u,

\]



где



\[

\begin{aligned}

& a(u)=\frac{1}{\pi} \int_{-\infty}^{\infty} f(t) \cos u t d t, \\

& b(u)=\frac{1}{\pi} \int_{-\infty}^{\infty} f(t) \sin u t d t .

\end{aligned}

\]



В частности, для четных функций



\[

f(x)=\int_{0}^{\infty} a(u) \cos u x d x

\]



где



\[

a(u)=\frac{2}{\pi} \int_{0}^{\infty} f(t) \cos u t d t .

\]



Формулу (2) можно рассматривать как предельную форму ряда Фурье для функций, имеюших период $2 T$, когда $T \rightarrow \infty$. При этом $a(u)$ и $b(u)$ аналогичны коэффициентам Фурье функции $f(x)$. Употребляя комплексные числа, можно за-





\end{document}